\begin{examplebox}
\textbf{例 1：同一支上的弦长，直接代入最快。}
直线 $y=x+2$ 与 $x^2/4-y^2=1$ 交于 $P,Q$，求横坐标和、积及弦长。

\textbf{解。} 图见前面的投影图。这里 $a^2=4,b^2=1$；直线写成 $x-y+2=0$，所以 $A=1,B=-1,C=2$。
先查 $M=-3\ne0$、$H=4-3=1>0$，确有两个交点。
本题数字简单，直接代入更短：
\[
 \frac{x^2}{4}-(x+2)^2=1
 \ \Longleftrightarrow\ 3x^2+16x+20=0.
\]
于是
\[
 x_1+x_2=-\frac{16}{3},\qquad x_1x_2=\frac{20}{3}.
\]
斜率为 $1$，故
\[
 PQ=\sqrt2\sqrt{\frac{256}{9}-\frac{80}{3}}
 =\boxed{\frac{4\sqrt2}{3}}.
\]
\textbf{回验。} 二次式恰可分解为 $(3x+10)(x+2)$，交点为
$(-10/3,-4/3)$、$(-2,0)$。二者均满足原方程，坐标差均为 $4/3$，距离吻合。

\textbf{快法的选择。} 熟悉公式时可直接用 $M,H$ 写出相同结果；数字容易展开时，一次代入往往更省力。公式用于减少计算，不必强行代替更短的计算。
\end{examplebox}
\clearpage
\begin{examplebox}
\textbf{例 2：纵开口，混合使用横、纵坐标关系。}
直线 $y=2x+1$ 与 $y^2/5-x^2/4=1$ 交于 $P,Q$，求 $x_1+x_2$、$y_1y_2$。
\begin{center}\begin{tikzpicture}[font=\small]
\draw[black!30,rounded corners=2pt] (0,0) rectangle (14.4,8.2);
\node[anchor=west,font=\small\bfseries] at (.5,7.72) {纵开口双曲线：两交点分居上下两支};
\begin{scope}[shift={(7.2,3.65)},scale=0.58]
\begin{scope}
\clip (-3.2,-4.1) rectangle (3.2,4.1);
\draw[blue,very thick,domain=-3.2:3.2,samples=161,variable=\t] plot (\t,{2.23606797749979*sqrt(1+(\t/2)^2)});
\draw[blue,very thick,domain=-3.2:3.2,samples=161,variable=\t] plot (\t,{-2.23606797749979*sqrt(1+(\t/2)^2)});
\end{scope}
\draw[-Stealth,black!45] (-3.2,0)--(3.2,0) node[right] {$x$};
\draw[-Stealth,black!45] (0,-4.1)--(0,4.1) node[above] {$y$};
\draw[orange!85!black,very thick] (-2.55,-4.1)--(1.55,4.1);\coordinate (P) at (-2.1356303077,-3.2712606154);\fill[black] (P) circle (1.8pt) node[above left=11pt] {$P$};
\coordinate (Q) at (0.6810848532,2.3621697063);\fill[black] (Q) circle (1.8pt) node[above left=12pt] {$Q$};
\end{scope}\node[anchor=west] at (10.4,5.7) {$y=2x+1$};\node[anchor=west] at (10.4,4.7) {$y_1y_2<0$};
\node[anchor=west] at (.7,.55) {$y^2/5-x^2/4=1$；先识别正号所在的轴，再选对应公式};
\end{tikzpicture}\end{center}

\textbf{解。} $a^2=5,b^2=4,A=2,B=-1,C=1$。
先算 $N=16-5=11\ne0$、$K=12>0$，再得到
\[
 x_1+x_2=-\frac{2\cdot2\cdot1\cdot4}{11}=\boxed{-\frac{16}{11}},
\qquad y_1y_2=-\frac{5(1+4\cdot4)}{11}=\boxed{-\frac{85}{11}}.
\]
\textbf{独立回验。} 代入原式得到 $11x^2+16x-16=0$，故 $x_1x_2=-16/11$。
于是 $y_1y_2=(2x_1+1)(2x_2+1)=4x_1x_2+2(x_1+x_2)+1=-85/11$，与公式一致。
\end{examplebox}
\clearpage
\begin{examplebox}
\textbf{例 3：已知弦长，反求截距。}
直线 $y=x+m$ 与 $x^2-2y^2=2$ 交于两个不同点，且 $PQ=4\sqrt3$，求 $m$。
\begin{center}\begin{tikzpicture}[font=\small]
\draw[black!30,rounded corners=2pt] (0,0) rectangle (14.4,8.2);
\node[anchor=west,font=\small\bfseries] at (.5,7.72) {两组答案互为中心对称：弦长相同};
\begin{scope}[shift={(7.2,3.65)},scale=0.5]
\begin{scope}
\clip (-7,-4.8) rectangle (7,4.8);
\draw[blue,very thick,domain=-4.8:4.8,samples=161,variable=\t] plot ({1.4142135623730951*sqrt(1+(\t/1)^2)},\t);
\draw[blue,very thick,domain=-4.8:4.8,samples=161,variable=\t] plot ({-1.4142135623730951*sqrt(1+(\t/1)^2)},\t);
\end{scope}
\draw[-Stealth,black!45] (-7,0)--(7,0) node[right] {$x$};
\draw[-Stealth,black!45] (0,-4.8)--(0,4.8) node[above] {$y$};
\draw[orange!85!black,very thick] (-6.8,-4.8)--(2.8,4.8);\draw[red!75!black,dashed,very thick] (-2.8,-4.8)--(6.8,4.8);\coordinate (P) at (-6.4494897428,-4.4494897428);\fill[black] (P) circle (1.8pt) node[above left=7pt] {$P$};
\coordinate (Q) at (-1.5505102572,0.4494897428);\fill[black] (Q) circle (1.8pt) ;
\coordinate (P') at (6.4494897428,4.4494897428);\fill[black] (P') circle (1.8pt) node[below right=7pt] {$P'$};
\coordinate (Q') at (1.5505102572,-0.4494897428);\fill[black] (Q') circle (1.8pt) ;
\node[anchor=south east] at (-2.8,.5) {$Q$};\node[anchor=north west] at (2.8,-.5) {$Q'$};\end{scope}\node[anchor=west,orange!85!black] at (11.5,6.5) {$m=2$};\node[anchor=west,red!75!black] at (11.1,4.5) {$m=-2$};
\node[anchor=west] at (.7,.55) {$x^2/2-y^2=1,\ y=x+m$；$(x,y)\mapsto(-x,-y)$ 使 $m\mapsto-m$};
\end{tikzpicture}\end{center}

\textbf{解。} 标准化为 $x^2/2-y^2=1$，有 $M=-1$、$H=m^2-1$。
两不同交点的必要且充分条件是 $|m|>1$。
由通用弦长公式直接得
\[
 PQ=\frac{2\sqrt2\sqrt{2(m^2-1)}}1=4\sqrt{m^2-1}.
\]
因此 $m^2-1=3$，得到 $\boxed{m=\pm2}$，均满足 $|m|>1$。
\textbf{回验。} 代入消元式 $x^2+4mx+2m^2+2=0$，判别式为 $8(m^2-1)=24$，
所以距离为 $\sqrt2\sqrt{24}=4\sqrt3$。中心对称使 $m=2$ 与 $m=-2$ 的弦等长，两个答案都应保留。
\end{examplebox}
\clearpage
% ===== 新增：用原例 1 的数字练习两种速推法（原例 1--3 保持不变）=====
\begin{examplebox}
\textbf{补例（复算例 1）：两点相减求和，交叉乘减求积。}
已知
\[
 \frac{x^2}{4}-y^2=1,\qquad y=x+2,
\]
两交点不同；这里 $a^2=4,b^2=1,k=1,m=2$，故
$\dfrac1{a^2}-\dfrac{k^2}{b^2}=\dfrac14-1=-\dfrac34\ne0$。

\textbf{求和：} 两点的双曲线方程直接相减，结合
$y_1-y_2=x_1-x_2$ 及 $y_1+y_2=(x_1+x_2)+4$，
得到
\[
 \frac{x_1+x_2}{4}=y_1+y_2=(x_1+x_2)+4,
 \qquad \boxed{x_1+x_2=-\frac{16}{3}}.
\]
\textbf{求积：} 两个代入方程交叉乘减后，
$2kmx_1x_2$ 自动抵消，留下
\[
 \left(\frac14-1\right)x_1x_2=-(1+2^2),
 \qquad \boxed{x_1x_2=\frac{20}{3}}.
\]
\textbf{核对：} 与原例 1 的 $3x^2+16x+20=0$ 的韦达结果完全相同，
但这次没有为求和、积而先展开整个消元方程。
\end{examplebox}
\clearpage

% ===== 再补充：原例 1 的完全逐步演算（此前例题原样保留）=====
\begin{examplebox}
\textbf{完整慢演示：同一道题，不先写二次方程，也能求和与积。}
\[
 \frac{x^2}{4}-y^2=1,\qquad y=x+2.
\]
设两个不同交点为 $P(x_1,y_1),Q(x_2,y_2)$。

\textbf{（一）求 $x_1+x_2$。} 先分别写出
\[
 \frac{x_1^2}{4}-y_1^2=1,\qquad
 \frac{x_2^2}{4}-y_2^2=1.
\]
两式相减、利用平方差，得到
\[
 \frac{(x_1-x_2)(x_1+x_2)}4
 =(y_1-y_2)(y_1+y_2).
\]
直线 $y=x+2$ 给出
\[
 y_1-y_2=x_1-x_2,\qquad y_1+y_2=x_1+x_2+4.
\]
因此直接约去非零的 $x_1-x_2$：
\[
 \frac{x_1+x_2}{4}=x_1+x_2+4.
\]
两边乘 $4$ 后整理：
\[
 x_1+x_2=4(x_1+x_2)+16,
 \qquad \boxed{x_1+x_2=-\frac{16}{3}}.
\]

\textbf{（二）求 $x_1x_2$。} 将 $y=x+2$ 放进双曲线，保留平方不急着整理：
\[
 F(x)=\frac{x^2}{4}-(x+2)^2-1=0.
\]
第一个交点给 $F(x_1)=0$，第二个给 $F(x_2)=0$。
第一个乘 $x_2$，第二个乘 $x_1$，相减：
\[
 \begin{aligned}
 0={}&x_2\left[\frac{x_1^2}{4}-(x_1+2)^2-1\right]\\
   &-x_1\left[\frac{x_2^2}{4}-(x_2+2)^2-1\right].
 \end{aligned}
\]
只展开必要部分：
\[
 (x_1+2)^2=x_1^2+4x_1+4,\quad
 (x_2+2)^2=x_2^2+4x_2+4.
\]
其中 $4x_1x_2-4x_2x_1=0$，交叉项全部抵消。
余下同含 $x_1-x_2$：
\[
 0=(x_1-x_2)\left[\left(\frac14-1\right)x_1x_2+5\right].
\]
因为 $x_1\ne x_2$：
\[
 -\frac34x_1x_2+5=0,
 \qquad \boxed{x_1x_2=\frac{20}{3}}.
\]
\textbf{回验（只用来检查答案）：} 若常规代入，会得到
$3x^2+16x+20=0$，韦达给出相同的和与积。
两点法的意义在于\textbf{即使记不起长公式，也能利用原方程快速恢复结论}。
\end{examplebox}
\clearpage
